\chapter{Buscher Rules: T-Duality from the Worldsheet}\label{ch:buscher}

In \cref{ch:tduality} T-duality appeared as a property of a spectrum: the mass formula of the
string on a circle does not change under $R\to\ap/R$, $n\leftrightarrow w$. That argument uses
the exact solution of a free theory. It says nothing about a string moving in a curved
background, where no mode expansion is available, and it does not explain \emph{why} the
two descriptions agree. This chapter answers a sharper question: given a sigma model with
metric $G$, Kalb--Ramond field $B$ and dilaton $\Phi$, and given only that the background
does not depend on one coordinate $\theta$, what is the dual background
$(\tilde G,\tilde B,\tilde\Phi)$? The answer is a set of algebraic formulas, the
\emph{Buscher rules}\index{Buscher rules}, and the derivation is a short manipulation of the
worldsheet path integral: make the shift symmetry of $\theta$ local, add a Lagrange multiplier
that undoes the gauging, and integrate out the fields in two different orders. One order
returns the original model; the other gives the dual one, with the Lagrange multiplier
promoted to the dual coordinate.

A student should care for three reasons. First, the method is the prototype for many
dualities in field theory: electric--magnetic duality in four dimensions and particle--vortex
duality in three can be derived in the same way (a standard fact, not checked against a source
in this book's library). Second, the rules are used every time T-duality
is applied to a background that is not a flat torus: D-brane solutions
(\cref{ch:dbranes}), backgrounds with flux (\cref{ch:tadpole}), and the torus fibres of
\cref{ch:k3t2}. Third, the rules look complicated in components but are one matrix
conjugation of the generalized metric $\HH(G,B)$; seeing this here prepares the $\Odd{d}$
action of \cref{ch:narain,ch:odd} and the doubled formalism of \cref{ch:doubled,ch:genmetric}.
The chapter closes with an account of what the Lean file
\lean{DoubleFieldTheory/TDualityBuscher.lean} states. It is an integer model of the
one-circle case, and we describe its scope precisely.

\paragraph{Conventions.} The worldsheet is Euclidean with coordinates $\sigma^a$, $a=1,2$, and
metric $h_{ab}$; $\varepsilon^{ab}$ is the antisymmetric symbol with $\varepsilon^{12}=+1$.
Target-space coordinates are $X^i$, $i=0,\dots,d-1$. As in the rest of the book $\ap$ has
dimension length$^2$ and $\ell_s=\sqrt{\ap}$. In this chapter the coordinate along the
isometry is called $\theta=X^0$ and is given the period
\begin{equation}\label{eq:buscher:period}
 \theta\sim\theta+2\pi\ell_s ,
\end{equation}
so that the metric components are dimensionless. A circle of physical radius $R$ is then
described by $G_{00}=R^2/\ap$, because the circumference is
$\sqrt{G_{00}}\cdot2\pi\ell_s=2\pi R$. This choice matters for the dilaton and for the global
analysis; \cref{sec:buscher:global} explains what changes for other periods.

%=====================================================================================
\section{The sigma model with an isometry}\label{sec:buscher:setup}

The starting point is the non-linear sigma model\index{sigma model} of \cref{ch:backgrounds},
written in Euclidean signature,
\begin{equation}\label{eq:buscher:sigma}
 S=\frac{1}{4\pi\ap}\int\dd^2\sigma\,\Big[\sqrt{h}\,h^{ab}G_{ij}(X)\,\partial_aX^i\partial_bX^j
   +\ii\,\varepsilon^{ab}B_{ij}(X)\,\partial_aX^i\partial_bX^j
   +\ap\sqrt{h}\,R^{(2)}\Phi(X)\Big],
\end{equation}
where $G_{ij}$ is the target-space metric, $B_{ij}=-B_{ji}$ the Kalb--Ramond
field\index{Kalb--Ramond field} (called ``torsion'' in part of the literature), $\Phi$ the
dilaton\index{dilaton} and $R^{(2)}$ the scalar curvature of $h_{ab}$
\source{AAL §2.1, ll.~217--235}. The factor $\ii$ in the $B$-term is a consequence of the
Wick rotation: a term with one time derivative, such as
$\varepsilon^{ab}\partial_aX^i\partial_bX^j$, picks up one factor of $\ii$ under
$\tau\to-\ii\tau$, whereas the kinetic term, with two, stays real. The $B$-term does not
involve $h_{ab}$ at all: it is the integral of the two-form $B$ pulled back to the worldsheet.

Assume that the background has an \emph{abelian isometry}\index{isometry}: there is a vector
field $k$ such that the action is invariant under $\delta X^i=\epsilon\,k^i(X)$. Invariance of
the three terms of \eqref{eq:buscher:sigma} requires $k$ to be a Killing vector of $G$, the
field strength $H=\dd B$ to be invariant, and $k^i\partial_i\Phi=0$
\source{GPR §4.1, ll.~5935--5953}. Under these conditions one can choose \emph{adapted
coordinates}\index{adapted coordinates} $X^i=(\theta,x^\alpha)$, $\alpha=1,\dots,d-1$, in which
$k=\partial_\theta$ and none of $G$, $B$, $\Phi$ depends on $\theta$; for $B$ this may require
a gauge transformation $B\to B+\dd\lambda$ \source{GPR §4.1, ll.~5954--5958}. In adapted
coordinates the isometry is the shift $\theta\to\theta+\epsilon$ and the Lagrangian of
\eqref{eq:buscher:sigma} splits as
\begin{align}\label{eq:buscher:split}
 4\pi\ap\,\mathcal L={}&\sqrt h\,h^{ab}\big(G_{00}\,\partial_a\theta\partial_b\theta
   +2G_{0\alpha}\,\partial_a\theta\,\partial_bx^\alpha
   +G_{\alpha\beta}\,\partial_ax^\alpha\partial_bx^\beta\big)\notag\\
 &+\ii\varepsilon^{ab}\big(2B_{0\alpha}\,\partial_a\theta\,\partial_bx^\alpha
   +B_{\alpha\beta}\,\partial_ax^\alpha\partial_bx^\beta\big)+\ap\sqrt h\,R^{(2)}\Phi .
\end{align}
The factor $2$ in the mixed $B$-term comes from the two orderings $(i,j)=(0,\alpha)$ and
$(\alpha,0)$: $B_{\alpha0}\varepsilon^{ab}\partial_ax^\alpha\partial_b\theta
=B_{0\alpha}\varepsilon^{ab}\partial_a\theta\partial_bx^\alpha$, since both $B$ and
$\varepsilon$ are antisymmetric. The field $\theta$ enters \eqref{eq:buscher:split} only
through its derivative $\partial_a\theta$. The whole construction depends on this fact.

For the derivations below we take the worldsheet metric flat, $h_{ab}=\delta_{ab}$, and write
$A\cdot B=A_aB_a$. Nothing is lost: in two dimensions one can always choose
$h_{ab}=\ee^{2\omega}\delta_{ab}$ locally, the kinetic term $\sqrt h h^{ab}$ is independent of
$\omega$, and the $\varepsilon$-terms never contained $h_{ab}$. The curvature term is the only
place where $h_{ab}$ matters, and it returns in \cref{sec:buscher:dilaton}.

%=====================================================================================
\section{Warm-up: the circle in three steps}\label{sec:buscher:circle}

Take first a single field $\theta$ with constant $G_{00}=R^2/\ap$ and nothing else:
\begin{equation}\label{eq:buscher:circleaction}
 S[\theta]=\frac{G_{00}}{4\pi\ap}\int\dd^2\sigma\;\partial_a\theta\,\partial_a\theta .
\end{equation}
With $\theta=\ell_s\varphi$ this is Tong's action $\frac{R^2}{4\pi\ap}\int(\partial\varphi)^2$
for a field $\varphi$ of period $2\pi$ \source{Tong §8.3.1, ll.~11718--11731}. We follow his
argument and keep track of the normalizations.

\paragraph{Step 1: gauge the shift symmetry.} The symmetry $\theta\to\theta+\lambda$ with
constant $\lambda$ is made local by introducing a worldsheet gauge field $A_a$ with
$A_a\to A_a-\partial_a\lambda$ and replacing $\partial_a\theta$ by the covariant derivative
$D_a\theta=\partial_a\theta+A_a$ \source{Tong §8.3.1, ll.~11732--11735}. This changes the
theory: $A_a$ is a new integration variable.

\paragraph{Step 2: add a Lagrange multiplier that undoes Step 1.} Introduce one more scalar
$\tilde\theta$ and the \emph{master action}\index{master action}\index{Lagrange multiplier}
\begin{equation}\label{eq:buscher:master1}
 S[\theta,A,\tilde\theta]=\frac{1}{4\pi\ap}\int\dd^2\sigma\,
   \Big[G_{00}\,D_a\theta\,D_a\theta+2\ii\,\tilde\theta\,\varepsilon^{ab}\partial_aA_b\Big].
\end{equation}
The new term is gauge invariant because
$\varepsilon^{ab}\partial_a\partial_b\lambda=0$. In terms of fields of period $2\pi$,
$\tilde\theta=\ell_s\vartheta$ and $A=\ell_s\mathcal A$, it reads
$\frac{\ii}{2\pi}\int\vartheta\,\varepsilon^{ab}\partial_a\mathcal A_b$, which is the
normalization of Tong's (8.12) \source{Tong §8.3.1, ll.~11736--11746}. The functional integral
over $\tilde\theta$ of $\exp(-S)$ is a functional delta function,
\[
 \int\mathcal D\tilde\theta\;
 \exp\Big(-\frac{\ii}{2\pi\ap}\int\dd^2\sigma\,\tilde\theta\,F\Big)\;\propto\;\delta[F],
 \qquad F=\varepsilon^{ab}\partial_aA_b=\partial_1A_2-\partial_2A_1 .
\]
This is why the multiplier term must be imaginary in Euclidean signature: a real exponent would
give a Gaussian damping, not a constraint. On a worldsheet with the topology of the plane or
the sphere, $F=0$ implies $A_a=\partial_a\chi$ for some function $\chi$; the gauge
transformation $\lambda=\chi$ then sets $A_a=0$ and \eqref{eq:buscher:master1} reduces to
\eqref{eq:buscher:circleaction}. On worldsheets with non-contractible cycles this conclusion
needs an extra argument, given in \cref{sec:buscher:global}.

\paragraph{Step 3: integrate in the other order.} Alternatively, use the gauge freedom first to
set $\theta=0$, so that $D_a\theta=A_a$, and integrate the multiplier term by parts,
\begin{equation}\label{eq:buscher:byparts}
 \int\dd^2\sigma\;\tilde\theta\,\varepsilon^{ab}\partial_aA_b
 =-\int\dd^2\sigma\;\varepsilon^{ab}\,\partial_a\tilde\theta\,A_b
 =\int\dd^2\sigma\;\varepsilon^{ab}A_a\,\partial_b\tilde\theta ,
\end{equation}
where the last step renames $a\leftrightarrow b$ and uses
$\varepsilon^{ba}=-\varepsilon^{ab}$. The action is now quadratic in $A_a$ with no
derivatives acting on it:
\begin{equation}\label{eq:buscher:quadratic1}
 4\pi\ap\,\mathcal L=G_{00}\,A_aA_a+2A_a\,j_a,\qquad
 j_a=\ii\,\varepsilon^{ab}\partial_b\tilde\theta .
\end{equation}
A Gaussian integral is evaluated by completing the square,
$G_{00}A\cdot A+2A\cdot j=G_{00}\,(A+j/G_{00})^2-j\cdot j/G_{00}$. The shifted integral gives
a field-independent factor (for constant $G_{00}$), and what remains is the value of the
action at the stationary point
\begin{equation}\label{eq:buscher:dualrel1}
 A_a=-\frac{\ii}{G_{00}}\,\varepsilon^{ab}\partial_b\tilde\theta .
\end{equation}
Since $\varepsilon^{ab}\varepsilon^{ac}=\delta^{bc}$, one has
$j\cdot j=-\partial_b\tilde\theta\,\partial_b\tilde\theta$. The two factors of $\ii$ turn
the sign of $-j\cdot j/G_{00}$ into a plus and produce a kinetic term of the correct sign:
\begin{equation}\label{eq:buscher:dualcircle}
 \tilde S[\tilde\theta]=\frac{1}{4\pi\ap}\,\frac{1}{G_{00}}\int\dd^2\sigma\;
   \partial_a\tilde\theta\,\partial_a\tilde\theta .
\end{equation}
This is again a circle sigma model, with $\tilde G_{00}=1/G_{00}=\ap/R^2$. If $\tilde\theta$
has the same period $2\pi\ell_s$ as $\theta$ (which \cref{sec:buscher:global} shows to be the
case), the dual radius is $\tilde R=\ell_s\sqrt{\tilde G_{00}}=\ap/R$
\source{Tong §8.3.1, ll.~11770--11799}. No mode expansion was used.

\begin{figure}[t]
\centering
\begin{tikzpicture}[>=Stealth,node distance=12mm,
  box/.style={draw,rounded corners=2pt,align=center,inner sep=4pt,font=\small}]
 \node[box] (m) at (0,0) {master action $S[\theta,A,\tilde\theta]$\\
    gauged isometry $+$ multiplier $\tilde\theta\,\varepsilon^{ab}\partial_aA_b$};
 \node[box] (o) at (-4.3,-2.6) {original model\\ $(G,B,\Phi)$, coordinate $\theta$};
 \node[box] (d) at (4.3,-2.6) {dual model\\ $(\tilde G,\tilde B,\tilde\Phi)$, coordinate
    $\tilde\theta$};
 \draw[->] (m.south west) -- node[left,align=right,font=\footnotesize,xshift=-2mm]
    {integrate $\tilde\theta$:\\ $F=0$, $A$ pure gauge} (o.north);
 \draw[->] (m.south east) -- node[right,align=left,font=\footnotesize,xshift=2mm]
    {fix $\theta=0$,\\ integrate $A$ (Gaussian)} (d.north);
 \draw[<->,dashed] (o) -- node[below,font=\footnotesize]
    {same partition function} (d);
\end{tikzpicture}
\caption{The logic of the Buscher procedure. Both sigma models are obtained from one path
integral by performing the integrations in different orders; they are therefore two
descriptions of the same quantum theory, provided the global conditions of
\cref{sec:buscher:global} hold.}
\label{fig:buscher:logic}
\end{figure}

\begin{physicsbox}[Momentum becomes winding]
Combine the two halves of the argument. Without fixing the gauge, the stationary point
\eqref{eq:buscher:dualrel1} reads $D_a\theta=-\frac{\ii}{G_{00}}\varepsilon^{ab}
\partial_b\tilde\theta$, because only the gauge-invariant combination
$D_a\theta=\partial_a\theta+A_a$ enters. The multiplier constraint allows the other gauge
choice, $A_a=0$, in which $D_a\theta=\partial_a\theta$. The two fields are therefore related
by
\begin{equation}\label{eq:buscher:dualrelcircle}
 G_{00}\,\partial_a\theta=-\ii\,\varepsilon^{ab}\partial_b\tilde\theta .
\end{equation}
The left-hand side is the Noether current of the shift symmetry of $\theta$; its charge is the
momentum along the circle. The right-hand side is a current that is conserved identically,
$\partial_a(\varepsilon^{ab}\partial_b\tilde\theta)=0$, whatever the equations of motion; its
charge counts how many times $\tilde\theta$ winds. So the duality maps the momentum of
$\theta$ to the winding of $\tilde\theta$, which is the exchange $n\leftrightarrow w$ of
\cref{ch:tduality}. After continuation to Lorentzian signature, where the factor $\ii$
disappears, \eqref{eq:buscher:dualrelcircle} is Tong's relation
$\partial_\alpha X=\epsilon_{\alpha\beta}\partial^\beta Y$ between the coordinate
$X=X_L+X_R$ and the dual coordinate $Y=X_L-X_R$ \source{Tong §8.3, ll.~11673--11689}.
\end{physicsbox}

%=====================================================================================
\section{A general background: deriving the Buscher rules}\label{sec:buscher:rules}

Return to the general Lagrangian \eqref{eq:buscher:split}. Since $\theta$ appears only through
$\partial_a\theta$, Steps 1--3 can be merged: after gauge fixing $\theta=0$ the covariant
derivative is just $A_a$, which we rename $V_a$. The master action is obtained from
\eqref{eq:buscher:split} by the substitution $\partial_a\theta\to V_a$, with $V_a$ an
independent field, plus the multiplier term \source{AAL §2.1, eq.~(2.1.2), ll.~236--248}:
\begin{align}\label{eq:buscher:master}
 4\pi\ap\,\mathcal L_{\rm master}={}&G_{00}\,V\!\cdot\! V+2G_{0\alpha}\,V\!\cdot\!\partial x^\alpha
  +G_{\alpha\beta}\,\partial x^\alpha\!\cdot\!\partial x^\beta\notag\\
 &+\ii\varepsilon^{ab}\big(2B_{0\alpha}V_a\partial_bx^\alpha
  +B_{\alpha\beta}\partial_ax^\alpha\partial_bx^\beta\big)
  +2\ii\,\tilde\theta\,\varepsilon^{ab}\partial_aV_b+\ap R^{(2)}\Phi .
\end{align}
Integrating over $\tilde\theta$ imposes $\varepsilon^{ab}\partial_aV_b=0$, hence
$V_a=\partial_a\theta$ on a topologically trivial worldsheet, and one is back at
\eqref{eq:buscher:split} \source{AAL §2.1, ll.~249--250}. The version with an explicit gauge
field and $\theta$ kept, due to Ro\v{c}ek and Verlinde, is equivalent and makes the geometric
meaning clearer: the dual model is obtained by gauging the isometry with a flat connection
\source{AAL §2.1, eq.~(2.1.8), ll.~342--370}.

Now integrate over $V_a$ instead. Use \eqref{eq:buscher:byparts} for the multiplier term and
collect everything linear in $V_a$:
\begin{equation}\label{eq:buscher:J}
 4\pi\ap\,\mathcal L_{\rm master}=G_{00}\,V\!\cdot\! V+2\,V\!\cdot\! J+(\text{terms without }V),
 \qquad
 J_a=G_{0\alpha}\partial_ax^\alpha+\ii\,\varepsilon^{ab}K_b ,
\end{equation}
with the abbreviation $K_b=\partial_b\tilde\theta+B_{0\alpha}\partial_bx^\alpha$. Completing
the square as before, the stationary point is
\begin{equation}\label{eq:buscher:Vsol}
 V_a=-\frac{J_a}{G_{00}}
 =-\frac{1}{G_{00}}\Big(G_{0\alpha}\partial_ax^\alpha
   +\ii\,\varepsilon^{ab}\big(B_{0\alpha}\partial_bx^\alpha+\partial_b\tilde\theta\big)\Big),
\end{equation}
in agreement with \source{AAL §2.1, eq.~(2.1.3), ll.~250--257}, and the $V$-dependent terms
are replaced by $-J\cdot J/G_{00}$. It remains to expand $J\cdot J$:
\begin{align}
 J\cdot J&=G_{0\alpha}G_{0\beta}\,\partial x^\alpha\!\cdot\!\partial x^\beta
  +2\ii\,G_{0\alpha}\,\varepsilon^{ab}\partial_ax^\alpha K_b
  -\varepsilon^{ab}\varepsilon^{ac}K_bK_c\notag\\
 &=G_{0\alpha}G_{0\beta}\,\partial x^\alpha\!\cdot\!\partial x^\beta
  -K\!\cdot\! K+2\ii\,G_{0\alpha}\,\varepsilon^{ab}\partial_ax^\alpha K_b ,
  \label{eq:buscher:JJ}\\
 K\!\cdot\! K&=\partial\tilde\theta\!\cdot\!\partial\tilde\theta
  +2B_{0\alpha}\,\partial\tilde\theta\!\cdot\!\partial x^\alpha
  +B_{0\alpha}B_{0\beta}\,\partial x^\alpha\!\cdot\!\partial x^\beta .\notag
\end{align}
Insert this into $-J\cdot J/G_{00}$ and sort the result by structure. The terms symmetric in
the worldsheet indices build the dual metric, those with $\varepsilon^{ab}$ build the dual
$B$-field:
\begin{itemize}[leftmargin=*,itemsep=1pt]
\item $\partial\tilde\theta\cdot\partial\tilde\theta$ has coefficient $1/G_{00}$;
\item $2\,\partial\tilde\theta\cdot\partial x^\alpha$ has coefficient $B_{0\alpha}/G_{00}$;
\item $\partial x^\alpha\cdot\partial x^\beta$ receives
 $-(G_{0\alpha}G_{0\beta}-B_{0\alpha}B_{0\beta})/G_{00}$ in addition to $G_{\alpha\beta}$;
\item the term $-\frac{2\ii}{G_{00}}G_{0\alpha}\varepsilon^{ab}\partial_ax^\alpha
 \partial_b\tilde\theta=+\frac{2\ii}{G_{00}}G_{0\alpha}\,\varepsilon^{ab}
 \partial_a\tilde\theta\,\partial_bx^\alpha$ has the form of the mixed $B$-term in
 \eqref{eq:buscher:split} with $B_{0\alpha}$ replaced by $G_{0\alpha}/G_{00}$;
\item the term $-\frac{2\ii}{G_{00}}G_{0\alpha}B_{0\beta}\,\varepsilon^{ab}
 \partial_ax^\alpha\partial_bx^\beta$ is antisymmetrized by $\varepsilon^{ab}$ and adds
 $-(G_{0\alpha}B_{0\beta}-G_{0\beta}B_{0\alpha})/G_{00}$ to $B_{\alpha\beta}$.
\end{itemize}
The outcome is a sigma model of the form \eqref{eq:buscher:split} for the coordinates
$(\tilde\theta,x^\alpha)$.

\begin{theorem}[Buscher rules]\label{thm:buscher:rules}
Let $(G,B,\Phi)$ be independent of $\theta=X^0$ and $G_{00}\neq0$. Integrating out $V_a$ in
\eqref{eq:buscher:master} gives the sigma model with background
\begin{equation}\label{eq:buscher:rules}
\begin{aligned}
 \tilde G_{00}&=\frac{1}{G_{00}}, &
 \tilde G_{0\alpha}&=\frac{B_{0\alpha}}{G_{00}}, &
 \tilde G_{\alpha\beta}&=G_{\alpha\beta}
   -\frac{G_{0\alpha}G_{0\beta}-B_{0\alpha}B_{0\beta}}{G_{00}},\\
 && \tilde B_{0\alpha}&=\frac{G_{0\alpha}}{G_{00}}, &
 \tilde B_{\alpha\beta}&=B_{\alpha\beta}
   -\frac{G_{0\alpha}B_{0\beta}-G_{0\beta}B_{0\alpha}}{G_{00}},
\end{aligned}
\end{equation}
together with the dilaton shift of \cref{sec:buscher:dilaton},
$\tilde\Phi=\Phi-\tfrac12\log G_{00}$. \TL\
\source{AAL §2.1, eqs.~(2.1.4)--(2.1.6), ll.~258--303; GPR §4.1, eq.~(4.1.9), ll.~6032--6065}
\end{theorem}

The extracted text of both sources garbles the layout of these equations. The assignment of
right-hand sides above is fixed by the derivation just given, and it agrees with GPR's
(4.1.9) once $B_{a0}=-B_{0a}$ is used there.

Three features deserve comment.

\emph{Metric and $B$-field mix.} The mixed components are exchanged,
$\tilde G_{0\alpha}\leftrightarrow B_{0\alpha}/G_{00}$ and
$\tilde B_{0\alpha}\leftrightarrow G_{0\alpha}/G_{00}$. In Kaluza--Klein language
(\cref{ch:circle}) $G_{0\alpha}/G_{00}$ is the gauge field that couples to momentum and
$B_{0\alpha}$ is the gauge field that couples to winding. The exchange of the two gauge fields
is the background-field version of $n\leftrightarrow w$.

\emph{A compact form.} In terms of $E_{ij}=G_{ij}+B_{ij}$ the five rules
\eqref{eq:buscher:rules} become four:
\begin{equation}\label{eq:buscher:Erules}
 \tilde E_{00}=\frac{1}{E_{00}},\qquad
 \tilde E_{0\alpha}=\frac{E_{0\alpha}}{E_{00}},\qquad
 \tilde E_{\alpha0}=-\frac{E_{\alpha0}}{E_{00}},\qquad
 \tilde E_{\alpha\beta}=E_{\alpha\beta}-\frac{E_{\alpha0}E_{0\beta}}{E_{00}} .
\end{equation}
To check, note $E_{00}=G_{00}$, and take symmetric and antisymmetric parts: for instance
$\tilde E_{0\alpha}+\tilde E_{\alpha0}=(E_{0\alpha}-E_{\alpha0})/E_{00}=2B_{0\alpha}/G_{00}
=2\tilde G_{0\alpha}$. The last rule is a Schur complement\index{Schur complement}: it is what
remains of a matrix after one row and column have been eliminated, as expected from a
Gaussian integration.

\emph{The rules are an involution.} The dual background is independent of $\tilde\theta$, so
the procedure can be repeated. Applying \eqref{eq:buscher:Erules} twice gives
$\tilde{\tilde E}_{00}=E_{00}$,
$\tilde{\tilde E}_{0\alpha}=(E_{0\alpha}/E_{00})\,E_{00}=E_{0\alpha}$,
$\tilde{\tilde E}_{\alpha0}=E_{\alpha0}$, and
$\tilde{\tilde E}_{\alpha\beta}=E_{\alpha\beta}-E_{\alpha0}E_{0\beta}/E_{00}
-(-E_{\alpha0}/E_{00})(E_{0\beta}/E_{00})E_{00}=E_{\alpha\beta}$. Dualizing twice returns the
original model \source{AAL §2.1, ll.~317--318}.

\begin{pitfallbox}[Signs and the orientation of the worldsheet]
The signs of the mixed components in \eqref{eq:buscher:rules}--\eqref{eq:buscher:Erules}
depend on conventions that differ between papers: the sign of the $B$-term in
\eqref{eq:buscher:sigma}, the sign of $\varepsilon^{12}$, and the sign of the multiplier
term. Flipping the orientation of the worldsheet, $\sigma^1\to-\sigma^1$, is equivalent to
$B\to-B$ \source{GPR §2.4, ll.~1423--1424} and changes $\tilde G_{0\alpha}$ and
$\tilde B_{0\alpha}$ by a sign, which can be absorbed by $\tilde\theta\to-\tilde\theta$. The
components $\tilde G_{00}$, $\tilde G_{\alpha\beta}$ and $\tilde B_{\alpha\beta}$ are quadratic
in the mixed fields and are not affected. Always check a formula taken from the literature on
one example before using it.
\end{pitfallbox}

%=====================================================================================
\section{The dilaton shift}\label{sec:buscher:dilaton}

The Gaussian integral over $V_a$ yields more than the stationary value of the action. It also
gives a determinant. At each worldsheet point the integral over the two components of $V_a$
contributes a factor proportional to $1/G_{00}(x(\sigma))$, so that, formally,
\begin{equation}\label{eq:buscher:formaldet}
 \int\mathcal DV\;\ee^{-S_{\rm master}}\;=\;\ee^{-\tilde S}\;\times
 \prod_{\sigma}\frac{1}{G_{00}\big(x(\sigma)\big)} .
\end{equation}
When $G_{00}$ depends on $x^\alpha$ this factor depends on the fields that remain and cannot be
dropped. An infinite product over worldsheet points has no meaning without regularization.
Buscher's result is that, with a regularization that respects conformal invariance, the
product equals the exponential of a local term of exactly the form of the dilaton coupling in
\eqref{eq:buscher:sigma}, and the net effect is\index{dilaton!shift under T-duality}
\begin{equation}\label{eq:buscher:dilaton}
 \tilde\Phi=\Phi-\frac12\log G_{00}
\end{equation}
\TL\ \source{AAL §2.1, eqs.~(2.1.6)--(2.1.7), ll.~294--316}. We give three arguments for
\eqref{eq:buscher:dilaton}, in increasing order of rigour.

\paragraph{(i) The low-energy observer.} For a circle of radius $R$ the effective action in
the non-compact dimensions is multiplied by $2\pi R\,\ee^{-2\Phi}=2\pi R/g_s^2$: the
circumference comes from the integral over the circle and $\ee^{-2\Phi}$ multiplies the
string-frame action of \cref{ch:backgrounds}. An observer who cannot distinguish $R$ from
$\tilde R=\ap/R$ must find the same lower-dimensional Newton constant in both descriptions,
$R/g_s^2=\tilde R/\tilde g_s^2$, hence
\begin{equation}\label{eq:buscher:gs}
 \tilde g_s=g_s\,\sqrt{\tilde R/R}=g_s\,\frac{\sqrt{\ap}}{R}
\end{equation}
\source{Tong §8.3, ll.~11691--11714}. With $g_s=\ee^{\Phi}$ and $G_{00}=R^2/\ap$ this is
$\tilde\Phi=\Phi-\frac12\log(R^2/\ap)$, i.e.\ \eqref{eq:buscher:dilaton}.

\paragraph{(ii) Counting on a lattice.} The next argument is a heuristic of our own, not taken
from the sources, and is valid for constant $G_{00}$. It shows where the Euler number comes
from. Replace the worldsheet by a cell decomposition with $N_0$ vertices, $N_1$ links and
$N_2$ faces, so that $N_0-N_1+N_2=\chi=2-2g$. The field $\theta$ lives on vertices, the
one-form $V$ on links, and the multiplier $\tilde\theta$, which multiplies the curl of $V$, on
faces. The reparametrization-invariant measure for $\theta$ is
$\prod_{\rm vertices}\sqrt{G_{00}}\,\dd\theta$, and for the dual field it is
$\prod_{\rm faces}\sqrt{\tilde G_{00}}\,\dd\tilde\theta
=\prod_{\rm faces}G_{00}^{-1/2}\,\dd\tilde\theta$. The Gaussian integral over $V$ gives
$G_{00}^{-1/2}$ per link. Starting from the original model with its invariant measure and
passing through the master action, the partition functions are related by
\[
 Z_{\rm orig}=G_{00}^{N_0/2}\cdot G_{00}^{-N_1/2}\cdot G_{00}^{N_2/2}\;Z_{\rm dual}
 =G_{00}^{\chi/2}\,Z_{\rm dual},
\]
where the last factor converts $\prod\dd\tilde\theta$ into the invariant dual measure. For
constant $\Phi$ the dilaton term in \eqref{eq:buscher:sigma} equals $\Phi\chi$, since
$\frac{1}{4\pi}\int\sqrt h R^{(2)}=\chi$, so each model carries a weight $\ee^{-\Phi\chi}$.
Equality of $\ee^{-\Phi\chi}Z_{\rm orig}$ and $\ee^{-\tilde\Phi\chi}Z_{\rm dual}$ requires
$\ee^{-\tilde\Phi\chi}=\ee^{-\Phi\chi}G_{00}^{\chi/2}$, which is \eqref{eq:buscher:dilaton}.
The local factors cancel between vertices, links and faces; only the topological remainder
survives, and a term proportional to $\chi$ is a dilaton.

\paragraph{(iii) The regularized determinant.} On the sphere, and in the gauged formulation,
one can write $A=\partial\alpha$, $\bar A=\bar\partial\beta$ in complex coordinates. The change of variables from
$(A,\bar A)$ to $(\alpha,\beta)$ has Jacobian $\det\Delta$, with
$\Delta=-\frac{1}{\sqrt h}\partial_a\sqrt h\,h^{ab}\partial_b$ the scalar Laplacian, and the
Gaussian integral over $\alpha,\beta$ produces $(\det\Delta_{G_{00}})^{-1}$ with
$\Delta_{G_{00}}=-\frac{1}{\sqrt h}\partial_a\sqrt h\,h^{ab}G_{00}\partial_b$. The measure
factor is thus the ratio $\det\Delta/\det\Delta_{G_{00}}$
\source{AAL §4, eqs.~(4.0.6)--(4.0.9), ll.~1013--1040}. Writing $G_{00}=1+\sigma$ and using
the heat-kernel expansion to first order in $\sigma$, Alvarez, Alvarez-Gaum\'e and Lozano
find
\begin{equation}\label{eq:buscher:heatkernel}
 \frac{\det\Delta_{G_{00}}}{\det\Delta}
  =\exp\Big(-\frac{1}{8\pi}\int\dd^2\sigma\sqrt h\,R^{(2)}\log G_{00}\Big)
\end{equation}
\TL\ \source{AAL §4, eqs.~(4.0.10)--(4.0.16), ll.~1048--1111}. The inverse of
\eqref{eq:buscher:heatkernel} multiplies $\ee^{-\tilde S}$. Compared with the dilaton term
$\exp\big(-\frac{1}{4\pi}\int\sqrt hR^{(2)}\Phi\big)$ it amounts to replacing $\Phi$ by
$\Phi-\frac12\log G_{00}$. The shift is needed for consistency and cannot be chosen
otherwise: without it the dual model is in general not conformally invariant, that is, its
beta functions (\cref{ch:backgrounds}) do not vanish even if those of the original model do
\source{AAL §4, ll.~945--953}. Buscher showed conformal invariance of the dual model with the
shift to first order in $\ap$; the gauging argument of Ro\v{c}ek and Verlinde extends the
equivalence to all orders \source{AAL §2.1, ll.~305--307 and 371--376}.

\begin{proposition}[An invariant measure]\label{prop:buscher:measure}
Under \eqref{eq:buscher:rules} and \eqref{eq:buscher:dilaton},
$\det\tilde G=\det G/G_{00}^{\,2}$ and therefore
\begin{equation}\label{eq:buscher:measure}
 \sqrt{\det\tilde G}\;\ee^{-2\tilde\Phi}=\sqrt{\det G}\;\ee^{-2\Phi}.
\end{equation}
\end{proposition}
\begin{proof}
Write $G$ in Kaluza--Klein form. With $a_\alpha=G_{0\alpha}/G_{00}$ one has
$G_{ij}\dd X^i\dd X^j=G_{00}(\dd\theta+a_\alpha\dd x^\alpha)^2
+g_{\alpha\beta}\dd x^\alpha\dd x^\beta$, where
$g_{\alpha\beta}=G_{\alpha\beta}-G_{0\alpha}G_{0\beta}/G_{00}$, and
$\det G=G_{00}\det g$ because the change of basis
$\dd\theta\to\dd\theta+a_\alpha\dd x^\alpha$ has unit determinant. For the dual metric
the same decomposition gives $\tilde g_{\alpha\beta}=\tilde G_{\alpha\beta}
-\tilde G_{0\alpha}\tilde G_{0\beta}/\tilde G_{00}
=\tilde G_{\alpha\beta}-B_{0\alpha}B_{0\beta}/G_{00}=g_{\alpha\beta}$. Hence
$\det\tilde G=\tilde G_{00}\det g=\det G/G_{00}^2$, and
$\sqrt{\det\tilde G}\,\ee^{-2\tilde\Phi}=(\sqrt{\det G}/G_{00})\,\ee^{-2\Phi}G_{00}$.
\end{proof}
The proof shows more than was claimed: the metric $g_{\alpha\beta}$ seen by a
lower-dimensional observer is duality invariant. The combination in
\eqref{eq:buscher:measure} is the measure of the string-frame effective action; in double
field theory it is written $\ee^{-2d}$ and $d$ is called the \emph{duality-invariant
dilaton}\index{dilaton!duality-invariant} (\cref{ch:dftaction}). For one circle,
$\sqrt{\det G}=R/\ell_s=\ee^{x}$ with $x=\log(R/\ell_s)$, so that
\begin{equation}\label{eq:buscher:twod}
 2d=2\Phi-x,\qquad x\mapsto-x,\qquad \Phi\mapsto\Phi-x .
\end{equation}
These three linear expressions are what the Lean file of \cref{sec:buscher:lean} encodes.

%=====================================================================================
\section{Global issues}\label{sec:buscher:global}

\subsection{The period of the Lagrange multiplier}

On a worldsheet of genus $g\ge1$ the constraint $\varepsilon^{ab}\partial_aV_b=0$, that is
$\dd V=0$, does not imply $V=\dd\theta$ for a single-valued function $\theta$. A closed
one-form can have non-zero \emph{holonomies}\index{holonomy}
$\nu_\gamma=\oint_\gamma V$ around the non-contractible cycles $\gamma$. For the original
model we need exactly the closed forms $V=\dd\theta$ with $\theta$ a map to the circle
\eqref{eq:buscher:period}; these have $\nu_\gamma\in2\pi\ell_s\Z$, the integer being the
number of times the worldsheet cycle winds around the target circle. Holonomies that are not
of this form would describe configurations that do not exist in the original model
(\cref{fig:buscher:torus}).

\begin{figure}[t]
\centering
\begin{tikzpicture}[>=Stealth,scale=0.95]
 \draw[thick] (0,0) rectangle (4,2.4);
 \draw[->,thick,blue!60!black] (0,0) -- (2.1,0);
 \draw[->,thick,blue!60!black] (0,2.4) -- (2.1,2.4);
 \draw[->,thick,red!60!black] (0,0) -- (0,1.3);
 \draw[->,thick,red!60!black] (4,0) -- (4,1.3);
 \node[below] at (2,0) {$a$};
 \node[left] at (0,1.2) {$b$};
 \node[font=\small,align=center] at (2,1.2) {$\dd V=0$\\[1mm]
   $\oint_aV=\nu_a,\ \oint_bV=\nu_b$};
 \node[font=\small,align=left,anchor=west] at (4.8,1.75)
   {$\oint_a\dd\tilde\theta=2\pi\tilde L\,m_a,\quad
     \oint_b\dd\tilde\theta=2\pi\tilde L\,m_b$};
 \node[font=\small,align=left,anchor=west] at (4.8,0.75)
   {$\displaystyle\sum_{m_a,m_b\in\Z}\Rightarrow\ \nu_a,\nu_b\in
     \frac{2\pi\ap}{\tilde L}\,\Z$};
\end{tikzpicture}
\caption{A torus worldsheet, drawn as a rectangle with opposite sides identified, and its two
cycles $a$ and $b$. The sum over the windings of the multiplier $\tilde\theta$ restricts the
holonomies of the flat gauge field $V$ to a lattice.}
\label{fig:buscher:torus}
\end{figure}

The multiplier enforces this if, and only if, it is itself periodic with the right period.
Let $\tilde\theta\sim\tilde\theta+2\pi\tilde L$. On the torus $\dd\tilde\theta$ is the sum
of an exact form and a harmonic form with periods $\oint_a\dd\tilde\theta=2\pi\tilde Lm_a$,
$\oint_b\dd\tilde\theta=2\pi\tilde Lm_b$, and the path integral over $\tilde\theta$ includes a
sum over the integers $m_a,m_b$. The multiplier term is defined by its integrated form
\eqref{eq:buscher:byparts}, $\frac{\ii}{2\pi\ap}\int V\wedge\dd\tilde\theta$. The exact
part of $\dd\tilde\theta$ imposes $\dd V=0$ as before. For two closed one-forms on the torus,
Riemann's bilinear identity gives
$\int V\wedge\dd\tilde\theta=\oint_aV\oint_b\dd\tilde\theta-\oint_bV\oint_a\dd\tilde\theta$,
so the harmonic part contributes
\begin{equation}\label{eq:buscher:holonomy}
 \sum_{m_a,m_b\in\Z}\exp\Big[-\frac{\ii\tilde L}{\ap}\big(\nu_a\,m_b-\nu_b\,m_a\big)\Big]
 \;\propto\;\sum_{k_a,k_b\in\Z}\delta\Big(\nu_a-\frac{2\pi\ap}{\tilde L}k_a\Big)\,
   \delta\Big(\nu_b-\frac{2\pi\ap}{\tilde L}k_b\Big),
\end{equation}
by the Poisson formula $\sum_m\ee^{-\ii my}=2\pi\sum_k\delta(y-2\pi k)$. The holonomies are
forced into $2\pi(\ap/\tilde L)\Z$. If $\theta$ has period $2\pi L$ we need them in
$2\pi L\,\Z$. The two periods are therefore tied together:\index{T-duality!global conditions}
\begin{equation}\label{eq:buscher:LL}
 L\,\tilde L=\ap .
\end{equation}
This is the content of the statements ``the holonomies are gauge trivial if $\theta$ has
periodicity $2\pi$'' \source{Tong §8.3.1, ll.~11759--11769} and ``this fixes the periodicity of
$\tilde\theta$ to be $2\pi$'' \source{GPR §4.1, ll.~6077--6126}, both in normalizations in which
$L=\tilde L$; the same analysis with Riemann's bilinear identity is in
\source{AAL §2.1, ll.~455--558}. With our choice $L=\ell_s$ one gets $\tilde L=\ell_s$, as
assumed in \cref{sec:buscher:circle}. Abelian isometries need only the torus analysis; higher
genus adds more pairs of cycles, each treated in the same way.

\begin{pitfallbox}[Where did the radius go?]
The rule $\tilde G_{00}=1/G_{00}$ is a statement about components in coordinates with
$L\tilde L=\ap$. If you prefer $X\sim X+2\pi R$ with $G_{XX}=1$, then $L=R$, the dual
coordinate has period $2\pi\ap/R$ and $\tilde G_{XX}=1$: the dual radius $\ap/R$ sits in the
period, not in the metric. For the dilaton the formula $\tilde\Phi=\Phi-\frac12\log G_{00}$ is
correct as written only when $L=\tilde L=\ell_s$; for general $L$ it reads
$\tilde\Phi=\Phi-\frac12\log(L^2G_{00}/\ap)$, as argument~(i) of
\cref{sec:buscher:dilaton} shows when the circumference is written as $2\pi L\sqrt{G_{00}}$.
A sign of trouble is a formula such as $\log G_{00}$ with a dimensionful $G_{00}$.
\end{pitfallbox}

For a non-compact isometry (a line instead of a circle) the periods of $\dd\tilde\theta$ are
arbitrary real numbers instead of multiples of $2\pi\tilde L$ \source{AAL §2.1,
ll.~509--510}. This is the limit $L\to\infty$, $\tilde L\to0$ of \eqref{eq:buscher:LL}. The sum
in \eqref{eq:buscher:holonomy} becomes an integral, which sets $\nu_a=\nu_b=0$: a field
$\theta$ valued in a line has no winding sectors, and the flat gauge field is pure gauge.

\subsection{Fixed points of the isometry}\label{sec:buscher:fixed}

The rules require $G_{00}\neq0$. At a fixed point of the isometry the Killing vector vanishes,
$G_{00}=k^ik_i=0$, and the stationary point \eqref{eq:buscher:Vsol} is singular
\source{AAL §2.1, ll.~331--332}. The simplest example is the plane in polar coordinates,
$\dd s^2=\dd r^2+r^2\dd\vartheta^2$, with the rotation $\vartheta\to\vartheta+\epsilon$ and
constant dilaton $\Phi_0$. In our conventions $\theta=\ell_s\vartheta$ and $G_{00}=r^2/\ap$,
and \cref{thm:buscher:rules} gives
\begin{equation}\label{eq:buscher:plane}
 \dd\tilde s^2=\dd r^2+\frac{\ap}{r^2}\,\dd\tilde\theta^2,\qquad
 \tilde\Phi=\Phi_0-\log\frac{r}{\ell_s}.
\end{equation}
A metric $\dd r^2+f(r)^2\dd\tilde\theta^2$ has scalar curvature $-2f''/f$; with
$f=\ell_s/r$ this is $-4/r^2$. The flat plane is mapped to a space whose curvature and
string coupling $\ee^{\tilde\Phi}=g_s\ell_s/r$ both diverge at $r=0$. Nothing is wrong with
the string on the plane. The dual sigma-model description breaks down where the orbit
of the isometry shrinks to zero size. A common interpretation (standard; not checked against
a source in this book's library) is that strings wound around the small orbit are light
there, and a description in terms of $\tilde G,\tilde B,\tilde\Phi$ alone does not contain
them. In such cases it is safer to work with the master action, which is regular
\source{AAL §2.1, ll.~331--332 and 574--576}.

\subsection{What the local derivation does not tell you}
Alvarez, Alvarez-Gaum\'e and Lozano list further limitations of the derivation in adapted
coordinates \source{AAL §2.1, ll.~319--334}. The choice of coordinates hides general covariance
and with it the global topology of the dual manifold: locally the dual is
$(M/S^1)\times S^1$, but whether the product is twisted is not visible; their example is the
round $S^3$, whose dual with respect to a fixed-point-free $U(1)$ is $S^2\times S^1$
\source{AAL §2.1, ll.~559--573}. Isometries that do not commute with the one being dualized
are in general not symmetries of the dual sigma model. The covariant formulation in terms of
the Killing vector $k$ replaces $G_{00}\to k^2$, $G_{0\alpha}\to k_\alpha$,
$B_{0\alpha}\to v_\alpha$, where $\iota_kH=-\dd v$ \source{AAL §2.1, eq.~(2.1.23),
ll.~578--604}.

%=====================================================================================
\section{Worked examples}\label{sec:buscher:examples}

\begin{example}[The circle, with numbers]\label{ex:buscher:circle}
Take $R=2\ell_s$ and $g_s=0.1$. Then $G_{00}=4$, $\tilde G_{00}=\frac14$, so
$\tilde R=\ell_s/2=\ap/R$. The dilaton shifts by $-\frac12\log4=-\log2$, so
$\tilde g_s=g_s/2=0.05$, in agreement with \eqref{eq:buscher:gs}. The invariant
\eqref{eq:buscher:measure} is $\sqrt{G_{00}}/g_s^2=2/0.01=200$ before and
$\frac12/0.0025=200$ after. In the variables of \eqref{eq:buscher:twod}, $x=\log2$ goes to
$-\log2$ and $2d=2\log0.1-\log2=\log0.005$ is unchanged. At the self-dual radius $R=\ell_s$
one has $G_{00}=1$, $x=0$, and neither the metric nor the dilaton changes.
\end{example}

\begin{example}[A tilted torus with $B$-field]\label{ex:buscher:numbers}
Let $d=2$ with constant background
\[
 G=\begin{pmatrix}2&1\\1&3\end{pmatrix},\qquad
 B=\begin{pmatrix}0&\frac32\\-\frac32&0\end{pmatrix},\qquad\det G=5 .
\]
Dualizing along $X^0$: $\tilde G_{00}=\frac12$, $\tilde G_{01}=B_{01}/G_{00}=\frac34$,
$\tilde B_{01}=G_{01}/G_{00}=\frac12$, and
$\tilde G_{11}=3-(1\cdot1-\frac32\cdot\frac32)/2=3+\frac58=\frac{29}{8}$. There is no
$\tilde B_{\alpha\beta}$ with $\alpha,\beta\ge1$ in two dimensions. Thus
\[
 \tilde G=\begin{pmatrix}\frac12&\frac34\\[2pt]\frac34&\frac{29}{8}\end{pmatrix},\qquad
 \tilde B=\begin{pmatrix}0&\frac12\\-\frac12&0\end{pmatrix},\qquad
 \det\tilde G=\tfrac{29}{16}-\tfrac{9}{16}=\tfrac54=\frac{\det G}{G_{00}^2},
\]
as \cref{prop:buscher:measure} requires, and $\tilde\Phi=\Phi-\frac12\log2$. Dualizing
again: $1/\tilde G_{00}=2$, $\tilde B_{01}/\tilde G_{00}=1$,
$\tilde G_{01}/\tilde G_{00}=\frac32$ and
$\frac{29}{8}-(\frac{9}{16}-\frac14)\cdot2=\frac{29}{8}-\frac58=3$: the original background
returns.
\end{example}

\begin{example}[$B$-field becomes tilt]\label{ex:buscher:tilt}
Let $G=\diag(G_{00},G_{11})$ be a rectangular two-torus with $B_{01}=b$. The rules give
\begin{equation}\label{eq:buscher:tilt}
 \tilde G=\begin{pmatrix}\dfrac{1}{G_{00}}&\dfrac{b}{G_{00}}\\[3mm]
   \dfrac{b}{G_{00}}&G_{11}+\dfrac{b^2}{G_{00}}\end{pmatrix},\qquad\tilde B=0,
 \qquad\text{i.e.}\qquad
 \dd\tilde s^2=\frac{1}{G_{00}}\big(\dd\tilde\theta+b\,\dd x^1\big)^2+G_{11}(\dd x^1)^2 .
\end{equation}
The dual torus carries no $B$-field, but it is no longer rectangular: the $B$-field has turned
into the angle between the two cycles. This is the simplest instance of the exchange of the
K\"ahler modulus $\rho$ (which contains $B$) and the complex-structure modulus $\tau$ (which
contains the angle) studied in \cref{ch:torus2}. The shift $b\to b+1$, which is a symmetry
of the original model because it changes the action by $2\pi\ii$ times an integer, becomes
the geometric transformation $\tilde\theta\to\tilde\theta-x^1$ of the dual torus.
\end{example}

%=====================================================================================
\section{Preview: one matrix instead of five rules}\label{sec:buscher:odd}

The rules \eqref{eq:buscher:rules} are non-linear in $(G,B)$. They become linear when $G$ and
$B$ are packaged into the \emph{generalized metric}\index{generalized metric} of
\cref{ch:genmetric}, the symmetric $2d\times2d$ matrix
\begin{equation}\label{eq:buscher:genmetric}
 \HH(G,B)=\begin{pmatrix}G-BG^{-1}B&BG^{-1}\\-G^{-1}B&G^{-1}\end{pmatrix}.
\end{equation}
Let $e=\diag(1,0,\dots,0)$ be the projector on the dualized direction and define
\begin{equation}\label{eq:buscher:Ttheta}
 T_\theta=\begin{pmatrix}1-e&-e\\-e&1-e\end{pmatrix},\qquad
 T_\theta^{\mathsf T}\,\eta\,T_\theta=\eta,\qquad T_\theta^2=1,\qquad
 \eta=\begin{pmatrix}0&1\\1&0\end{pmatrix}.
\end{equation}

\begin{proposition}\label{prop:buscher:conjugation}
With $(\tilde G,\tilde B)$ given by \eqref{eq:buscher:rules},
\begin{equation}\label{eq:buscher:conjugation}
 \HH(\tilde G,\tilde B)=T_\theta^{\mathsf T}\;\HH(G,B)\;T_\theta .
\end{equation}
\end{proposition}
We have verified \eqref{eq:buscher:conjugation} for $d=2$ and $d=3$ with fully symbolic $G$
and $B$ using sympy; this is a symbolic check, not a Lean theorem. The proof for all $d$ is a
block-matrix computation (\cref{exr:buscher:conj}); the general $\Odd{d}$ action is the
subject of \cref{ch:odd}. For $d=1$ it is immediate:
$\HH=\diag(G_{00},1/G_{00})$, $T_\theta=-\eta$, and conjugation swaps the two entries. The
matrix $T_\theta$ has integer entries and preserves $\eta$, so it belongs to $\Odd{d}$. It
is, up to signs, the \emph{factorized duality}\index{factorized duality} $D_0$ of Giveon,
Porrati and Rabinovici, whose off-diagonal blocks are $+e$ \source{GPR §2.4, eq.~(2.4.29),
ll.~1397--1421}. The two matrices are related by $B\to-B$, the orientation convention of the
pitfall box in \cref{sec:buscher:rules}: with $\HH(G,-B)$ on both sides,
\eqref{eq:buscher:conjugation} holds with $+e$ (also checked symbolically). GPR state that
the sigma-model duality acts on the background ``exactly as in the flat case''
\source{GPR §4.1, ll.~6025--6031}. In the same way \eqref{eq:buscher:Erules} is the
fractional-linear map $\tilde E=\big((1-e)E-e\big)\big({-e}E+(1-e)\big)^{-1}$.

Three consequences follow. Dualities along different circles are conjugations by
commuting matrices; dualizing all $d$ directions with $B=0$ gives
$\eta\,\HH(G,0)\,\eta=\HH(G^{-1},0)$, the matrix version of $R\to\ap/R$; and the
quantity \eqref{eq:buscher:measure} is the only place where the dilaton enters, which is
why double field theory uses $(\HH,d)$ as its field variables (\cref{ch:dftaction}).

%=====================================================================================
\section{What the machine has checked}\label{sec:buscher:lean}

The file \lean{DoubleFieldTheory/TDualityBuscher.lean} belongs to one of the older libraries
of the project. It contains no sigma model, no path integral and no metric with more than one
component. It models the one-circle case with integers and proves arithmetic identities. We
quote every theorem of the file and say what each one states.

\subsection{The matrix part}
The imported module \lean{DoubleFieldTheory.GeneralizedGeometry} defines a structure
\lean{Mat2} of $2\times2$ integer matrices with fields \lean{a b c d}, the product
\lean{MatMul}, the transpose \lean{MatTranspose}, the matrix
\lean{InversionGen}${}=\left(\begin{smallmatrix}0&1\\1&0\end{smallmatrix}\right)$ and the
diagonal matrix \lean{GenMetric g ginv}${}=\diag(g,g_{\rm inv})$ with two \emph{independent}
integer arguments.

\begin{leanbox}[Congruence action and the self-dual point]
\begin{leancode}
def CongruenceAction (M H : Mat2) : Mat2 :=
  MatMul (MatTranspose M) (MatMul H M)

theorem t_duality_congruence (H : Mat2) (hH : H = GenMetric 1 1) :
    CongruenceAction InversionGen H = GenMetric 1 1 := by ...

theorem inversion_generator_properties :
    MatMul InversionGen InversionGen = Identity2 ∧ Det2 InversionGen = -1 := by
  decide
\end{leancode}
\textbf{Reading.} \lean{CongruenceAction} is $M^{\mathsf T}HM$, the $d=1$ form of
\eqref{eq:buscher:conjugation}. The first theorem says: if $H$ is the $2\times2$ identity
matrix then $\sigma^{\mathsf T}H\sigma$ is the identity matrix, where
$\sigma=\eta$ for $d=1$. It concerns one fixed matrix, the generalized metric at the self-dual
radius $G_{00}=1$; it does not state $\diag(g,g^{-1})\mapsto\diag(g^{-1},g)$ for other radii.
The second theorem says $\sigma^2=1$ and $\det\sigma=-1$.
\textbf{Proof idea.} After substituting the hypothesis, both sides of the first statement
reduce by computation to the same structure and \lean{rfl} closes the goal; the second is a
closed statement about integers settled by \lean{decide}. \TA
\end{leanbox}

\subsection{The logarithmic radius and the dilaton}
The remaining declarations use an integer $x$ in the role of $\log(R/\ell_s)$ and an integer
\lean{phi} in the role of $\Phi$, following \eqref{eq:buscher:twod}.

\begin{leanbox}[Four integer identities]
\begin{leancode}
def BuscherLogMap (x : Int) : Int := -x
def BuscherDilatonMap (phi x : Int) : Int := phi - x
def TwoDilaton (phi x : Int) : Int := 2 * phi - x

theorem buscher_log_involution (x : Int) :
    BuscherLogMap (BuscherLogMap x) = x := by ...

theorem buscher_dilaton_involution (phi x : Int) :
    BuscherDilatonMap (BuscherDilatonMap phi x) (BuscherLogMap x) = phi :=
  by ...

theorem buscher_dilaton_measure_invariance (phi x : Int) :
    TwoDilaton (phi - x) (-x) = TwoDilaton phi x := by ...

theorem self_dual_radius_rigidity (x : Int) (hx : BuscherLogMap x = x) :
    x = 0 := by ...
\end{leancode}
\textbf{Reading.} For all integers: $-(-x)=x$; $(\varphi-x)-(-x)=\varphi$;
$2(\varphi-x)-(-x)=2\varphi-x$; and $-x=x$ implies $x=0$. These are the \emph{arithmetic
shadows} of four statements of this chapter: the map $R\to\ap/R$ is an involution; the dilaton
shift \eqref{eq:buscher:dilaton} is undone by a second dualization (because the second shift is
computed with the dual radius); the combination $2d$ of \eqref{eq:buscher:twod} is invariant;
and the self-dual radius is the only fixed point.
\textbf{Proof idea.} Each proof unfolds the definitions with \lean{dsimp} and calls
\lean{omega}, the decision procedure for linear integer arithmetic. \TA\ for the integer
identities.
\end{leanbox}

\paragraph{Honest scope.} An integer cannot represent $\log(R/\ell_s)$ for a general radius,
so these theorems are not statements about the real numbers $x$ and $\Phi$. The four
identities are linear and hold verbatim over $\R$ (a one-line \lean{ring} or \lean{linarith}
proof), but that transfer is not in the file. The documentation strings in the file describe
the physics in stronger terms than the statements justify; when the two differ, read the
statement. The identification of \lean{x} with a logarithmic radius and of \lean{phi} with
the dilaton is an interpretation and is never \TA. None of the following is formalized in
this file, nor, as far as a search of the catalogue of compiled declarations shows, elsewhere
in the project: the master action, the Gaussian integration, the rules
\eqref{eq:buscher:rules} for $d>1$, the determinant \eqref{eq:buscher:heatkernel}, and the
periodicity argument \eqref{eq:buscher:LL}.

\paragraph{What is proved elsewhere.} The newer library \lean{DualScaleStream2} works with
Mathlib matrices. Its file \lean{TDuality/Factorized.lean} defines the factorized duality for
any $d$, with GPR's signs, and proves that it is an involution and an element of $\Odd{d}$.
Its file \lean{DFT/GeneralizedMetric.lean} proves that the full duality inverts the metric at
$B=0$:
\begin{leancode}
def factorized (k : Fin d) : Matrix (Charge d) (Charge d) ℤ :=
  fromBlocks (1 - proj k) (proj k) (proj k) (1 - proj k)

theorem factorized_mul_self (k : Fin d) :
    factorized k * factorized k = 1 := by ...
theorem factorized_isODD (k : Fin d) : IsODD (factorized k) := by ...

theorem tduality_inverts_metric (G : Matrix (Fin d) (Fin d) ℝ)
    (hG : IsUnit G.det) :
    etaR d * genMetric G 0 * etaR d = genMetric G⁻¹ 0 := by ...
\end{leancode}
These are statements about real and integer matrices for all $d$; they are discussed in
\cref{ch:odd,ch:genmetric}. A faithful formalization of \cref{thm:buscher:rules} as an
algebraic statement would define the map $(G,B)\mapsto(\tilde G,\tilde B)$ on real matrices
with $G_{00}\neq0$ and prove \cref{prop:buscher:conjugation} and the involution property.
This is within reach of Mathlib's block-matrix and Schur-complement lemmas and is a natural
project for a reader (see also \cref{ch:open}). A formalization of the path-integral derivation is not, since Mathlib has
no functional integral.

\begin{center}
\small
\begin{tabular}{@{}>{\raggedright\arraybackslash}p{0.43\textwidth}c
  >{\raggedright\arraybackslash}p{0.45\textwidth}@{}}
\toprule
Claim & Tier & Evidence\\
\midrule
$\sigma^{\mathsf T}\mathbf 1\sigma=\mathbf 1$, $\sigma^2=1$, $\det\sigma=-1$ (integer
 $2\times2$) & \TA & \lean{t\_duality\_congruence}, \lean{inversion\_generator\_properties}\\[2pt]
$-(-x)=x$; $-x=x\Rightarrow x=0$ over $\Z$ & \TA & \lean{buscher\_log\_involution},
 \lean{self\_dual\_radius\_rigidity}\\
$(\varphi-x)-(-x)=\varphi$; $2(\varphi-x)+x=2\varphi-x$ over $\Z$ & \TA &
 \lean{buscher\_dilaton\_involution}, \lean{buscher\_dilaton\_measure\_invariance}\\
$D_k^2=1$, $D_k\in\Odd{d}$; $\eta\HH(G,0)\eta=\HH(G^{-1},0)$ & \TA &
 \lean{factorized\_mul\_self}, \lean{factorized\_isODD},
 \lean{tduality\_inverts\_metric}\\
Buscher rules \eqref{eq:buscher:rules} from the master action & \TL & AAL §2.1; GPR §4.1;
 derived on the page\\
Dilaton shift \eqref{eq:buscher:dilaton} from the regularized determinant & \TL & AAL §4\\
Equivalence on all genera given $L\tilde L=\ap$ & \TL & AAL §2.1; GPR §4.1\\
$\HH(\tilde G,\tilde B)=T_\theta^{\mathsf T}\HH T_\theta$ for $d=2,3$ & --- & symbolic check
 (sympy)\\
\bottomrule
\end{tabular}
\end{center}

%=====================================================================================
\begin{summarybox}
\begin{itemize}[leftmargin=*,itemsep=1pt]
\item If $G$, $B$, $\Phi$ do not depend on a coordinate $\theta$, the sigma model can be
 rewritten with an auxiliary vector $V_a$ in place of $\partial_a\theta$ and a Lagrange
 multiplier $\tilde\theta$ imposing $\varepsilon^{ab}\partial_aV_b=0$.
\item Integrating out $\tilde\theta$ returns the original model. Integrating out $V_a$, a
 Gaussian integral, gives a sigma model for $(\tilde\theta,x^\alpha)$ with the Buscher
 background \eqref{eq:buscher:rules}: $\tilde G_{00}=1/G_{00}$, and the mixed components of
 $G$ and $B$ are exchanged.
\item The determinant of the Gaussian integral shifts the dilaton,
 $\tilde\Phi=\Phi-\frac12\log G_{00}$; then $\sqrt{\det G}\,\ee^{-2\Phi}$ is invariant.
\item On worldsheets of higher genus the equivalence requires the periods of $\theta$ and
 $\tilde\theta$ to satisfy $L\tilde L=\ap$. The rules fail where the Killing vector
 vanishes.
\item In terms of the generalized metric the rules are one conjugation by an element of
 $\Odd{d}$.
\item In Lean, the file for this chapter proves integer identities that mirror the one-circle
 case ($x\to-x$, $\Phi\to\Phi-x$, $2d$ invariant, unique fixed point $x=0$). The rules
 themselves are Tier~L.
\end{itemize}
\end{summarybox}

\section*{Exercises}
\addcontentsline{toc}{section}{Exercises}

\begin{exercise}[$\star$]\label{exr:buscher:gauss}
Evaluate the ordinary integral $\int\dd^2V\exp\big(-\frac{1}{4\pi\ap}(gV\cdot V+2V\cdot j)\big)$
for a constant $g>0$ and a complex vector $j_a$. Identify the stationary point, the exponent
and the prefactor, and explain which of the three becomes the dilaton shift.
\end{exercise}

\begin{exercise}[$\star$]\label{exr:buscher:numbers}
Dualize the background of \cref{ex:buscher:numbers} along $X^1$ instead of $X^0$ (relabel the
coordinates, apply \eqref{eq:buscher:rules}, relabel back). Check
$\det\tilde G=\det G/G_{11}^2$.
\end{exercise}

\begin{exercise}[$\star$]\label{exr:buscher:Eform}
Derive \eqref{eq:buscher:Erules} from \eqref{eq:buscher:rules}, and conversely.
\end{exercise}

\begin{exercise}[$\star\star$]\label{exr:buscher:complex}
Repeat the derivation of \cref{sec:buscher:rules} in complex worldsheet coordinates, starting
from GPR's action $\frac{1}{2\pi}\int\dd^2z\,E_{ij}\partial X^i\bar\partial X^j$ with gauge
fields $A,\bar A$ and multiplier term $\tilde\theta(\partial\bar A-\bar\partial A)$
\source{GPR §4.1, ll.~5976--6000}. Show that the integration over $A,\bar A$ leads directly
to \eqref{eq:buscher:Erules}, up to the sign conventions discussed in the pitfall box.
\end{exercise}

\begin{exercise}[$\star\star$]\label{exr:buscher:gprdilaton}
GPR write the dilaton coupling as $-\frac{1}{8\pi}\int\dd^2z\,\phi R^{(2)}$ and the shift as
$\phi'=\phi+\log G_{00}$ \source{GPR §4.1, ll.~5947--5953 and 6067--6071}. Find the relation
between their $\phi$ and the $\Phi$ of this chapter that makes the two statements agree.
(Hint: compare the coefficient of the curvature term and apply both shifts to the circle.)
\end{exercise}

\begin{exercise}[$\star\star$]\label{exr:buscher:hopf}
The round three-sphere of radius $r_0$ can be written as
$\dd s^2=\frac{r_0^2}{4}\big[\dd\vartheta^2+\sin^2\!\vartheta\,\dd\phi^2
+(\dd\psi+\cos\vartheta\,\dd\phi)^2\big]$, with $\psi\sim\psi+4\pi$. Set $B=0$ and dualize
along $\psi$ (first rescale $\psi$ so that its period is $2\pi\ell_s$). Show that the dual
metric is the product of a round $S^2$ and a circle, with
$\tilde B\propto\cos\vartheta\,\dd\tilde\theta\wedge\dd\phi$, and that the dilaton is
constant. Compare with \source{AAL §2.1, ll.~570--573}.
\end{exercise}

\begin{exercise}[$\star\star$, Lean]\label{exr:buscher:leangeneral}
State and prove in Lean the $d=1$ Buscher rule that is missing from the file:
\begin{leancode}
theorem buscher_circle_shadow (g ginv : Int) :
    CongruenceAction InversionGen (GenMetric g ginv) = GenMetric ginv g
\end{leancode}
(Hint: \lean{simp [CongruenceAction, MatMul, MatTranspose, InversionGen, GenMetric]}.) Why
does this statement not need the hypothesis $g\cdot g_{\rm inv}=1$?
\end{exercise}

\begin{exercise}[$\star\star$, Lean]\label{exr:buscher:leanreal}
Using Mathlib, restate the last two theorems of the second box of
\cref{sec:buscher:lean} (the invariance of \lean{TwoDilaton} and the uniqueness of the fixed
point) for real \lean{x} and \lean{phi}, and prove them. Which tactic replaces \lean{omega}, and why
does \lean{omega} not apply?
\end{exercise}

\begin{exercise}[$\star\star\star$]\label{exr:buscher:conj}
Prove \cref{prop:buscher:conjugation} for all $d$. (Hint: write
$\HH=\left(\begin{smallmatrix}1&B\\0&1\end{smallmatrix}\right)
\left(\begin{smallmatrix}G&0\\0&G^{-1}\end{smallmatrix}\right)
\left(\begin{smallmatrix}1&0\\-B&1\end{smallmatrix}\right)$, and use the Kaluza--Klein
decomposition of the proof of \cref{prop:buscher:measure} to reduce the statement to $d=1$
blocks.)
\end{exercise}

\section*{Further reading}
\addcontentsline{toc}{section}{Further reading}
\begin{itemize}[leftmargin=*,itemsep=1pt]
\item The path-integral argument for one circle, including the remark on holonomies:
 \source{Tong §8.3.1, ll.~11716--11803}; the dilaton from the effective action:
 \source{Tong §8.3, ll.~11691--11714}.
\item Buscher's formulation, the gauging of Ro\v{c}ek and Verlinde, the covariant form of the
 rules and the global analysis: \source{AAL §2.1, ll.~217--604}. The dilaton and the
 regularized determinant: \source{AAL §4, ll.~943--1111}.
\item Factorized duality in sigma models in complex coordinates and the periodicity of the
 multiplier: \source{GPR §4.1, ll.~5904--6127}; the matrix $D_i$ in the flat case:
 \source{GPR §2.4, ll.~1397--1424}.
\item Buscher's original papers are cited in AAL as reference [16] and in GPR as references
 [33]--[36]; they are not part of this book's library and have not been checked directly.
\end{itemize}
